\subsection{Model Experimental Design}

The modeling stage consumes the processed development and test files produced by
the upstream data pipeline. For each of the three datasets, the data pipeline
creates a stratified 80/20 development--test split. The test partition is held
out during model comparison and is used only after cross-validation has been
completed. All three candidate algorithms receive the same processed feature
matrix for a given dataset so that differences in results are attributable to
the model rather than to different input representations.

Within the development partition, we apply stratified five-fold
cross-validation with shuffling and a fixed random seed of 42. Stratification
preserves approximately the same class proportion in each fold, which is
particularly important when the positive healthcare outcome is less frequent
\cite{pedregosa2011scikit}. The exact same fold-generation rule and random seed
are used for AdaBoost, XGBoost, and LightGBM. For every fold, four folds are used
for model fitting and the remaining fold is used for validation. Model objects
are re-created from scratch for every fold to prevent fitted state from leaking
between folds.

To reduce the effect of class imbalance during fitting, balanced sample weights
are computed from the training portion of each fold only. Validation and test
sets remain untouched so that the reported metrics reflect the natural class
prevalence of each dataset. A fixed probability threshold of 0.50 is used for
the threshold-dependent metrics. No threshold is optimized on the validation
folds, which avoids adding a second tuning procedure that could make fold
comparisons less consistent.

\subsection{Model Configuration and Hyperparameters}

The core comparison uses one fixed, literature-guided configuration per
algorithm across all three datasets. We intentionally do not select a different
best hyperparameter configuration for each dataset using the same five folds
that are later used to quantify stability. Doing so could make the reported
cross-validation scores optimistic and would confound RQ2 with hyperparameter
search variability. The selected settings use conservative learning rates,
moderate ensemble sizes, and regularization/subsampling for the gradient-
boosted models.

\begin{table}[h]
\centering
\small
\begin{tabular}{p{2.3cm}p{5.2cm}p{7.1cm}}
\textbf{Model} & \textbf{Selected hyperparameters} & \textbf{Rationale} \\
\hline
AdaBoost &
Decision stump (max depth = 1); 200 estimators; learning rate = 0.10 &
Decision stumps are the standard weak learner for AdaBoost. The larger ensemble
combined with a smaller learning rate provides gradual boosting while limiting
the contribution of each individual stump. \\
\hline
XGBoost &
300 estimators; learning rate = 0.05; max depth = 4; subsample = 0.80;
column sample = 0.80; $L_2$ regularization = 1; histogram tree method &
A low learning rate and moderate tree depth control model complexity. Row and
column subsampling reduce variance, while $L_2$ regularization makes leaf
weights more conservative. The histogram method improves efficiency on the
larger datasets. \\
\hline
LightGBM &
300 estimators; learning rate = 0.05; 31 leaves; max depth = 6; minimum child
samples = 20; subsample = 0.80; column sample = 0.80; $L_2$ regularization = 1 &
The number of leaves and depth cap constrain the leaf-wise tree growth. The
minimum child size, subsampling, and regularization reduce overfitting while the
small learning rate supports a smoother boosting process. Deterministic CPU
settings and a fixed seed are enabled for reproducibility. \\
\end{tabular}
\caption{Fixed model configurations used for the cross-dataset comparison.}
\label{tab:model_hyperparameters}
\end{table}

\subsection{Evaluation Metrics and Model Comparison}

Each validation fold is evaluated using ROC-AUC, PR-AUC, Recall, and F1-score.
ROC-AUC measures ranking performance across classification thresholds. PR-AUC
focuses on the precision--recall trade-off and is especially informative when
the positive class is relatively rare \cite{saito2015precisionrecall}. Recall
measures the proportion of positive cases correctly identified, while F1-score
balances precision and recall at the fixed 0.50 threshold. Because the three
datasets may have different class prevalence, PR-AUC is used as the primary
ranking metric, with ROC-AUC, F1-score, and fold variability retained for a more
complete comparison.

For RQ1, we report the mean five-fold score of every metric for every
model--dataset pair and rank the three models within each dataset by mean
PR-AUC. We additionally report the average within-dataset rank across all three
datasets rather than pooling patient rows across datasets. This prevents the
largest dataset from dominating the overall comparison.

For RQ2, predictive stability is measured directly from the five fold-level
scores. For each metric we report the mean, standard deviation, minimum,
maximum, and range across folds. A smaller standard deviation and range indicate
that the model is less sensitive to the particular development-data partition.
After cross-validation, each model is re-fitted once on the full development
partition and evaluated on the untouched test partition. These test results are
reported separately from the cross-validation results and are not used to
select the best model.

The complete experiment configuration is stored in \texttt{config.yaml}. The
implementation writes fold-level metrics, aggregated cross-validation summaries,
held-out test metrics, per-dataset rankings, cross-dataset rankings, and an
experiment metadata file containing package versions and a hash of the
configuration. This makes the model experiments repeatable and provides the
fold-level outputs required by the later SHAP stability analysis.
